\section{Hệ tọa độ, dấu góc và điều kiện quan sát}
\subsection{Quy ước trục và phép thử nghiêng từng trục}
Mọi phép tính góc chỉ có ý nghĩa khi quy ước trục được giữ nhất quán từ cảm biến, chương trình đến đồ thị. Trong báo cáo này, $X,Y,Z$ là ba trục gắn với module MPU6500. Roll mô tả phép nghiêng chủ yếu quanh $X$, Pitch quanh $Y$. Khi đặt board phẳng theo tư thế hiệu chuẩn, thành phần gia tốc trọng trường có độ lớn gần $1g$ trên trục $Z$. Dấu của Roll và Pitch cần được xác nhận bằng cách nghiêng board theo một chiều đã chọn, xem đồng thời dấu của gia tốc, gyro và góc hiển thị. Nếu lắp module quay $90\degree$ so với trục của bàn servo, việc so sánh trực tiếp Roll với góc lệnh sẽ sai dù thuật toán chạy đúng.

Phép thử một trục có một ưu thế quan trọng: có thể so sánh định tính bốn phương pháp trên cùng kích thích và cùng hệ trục. Tuy nhiên, khi quay đồng thời nhiều trục, góc Euler phụ thuộc thứ tự quay. Các phép xấp xỉ tích phân trực tiếp $\omega_x$ thành Roll và $\omega_y$ thành Pitch khi đó tạo sai số hình học. Để nghiên cứu chuyển động ba chiều đầy đủ phải dùng mô hình động học Euler với hiệu chỉnh chéo hoặc quaternion. Đề tài K03 tập trung vào phép nghiêng một trục và hai góc Roll/Pitch; mọi kết luận về độ chính xác được giới hạn trong phạm vi này.

\subsection{Tại sao gia tốc kế không đo được Yaw}
Ở trạng thái tĩnh, gia tốc kế quan sát hướng của vector trọng lực trong hệ trục cảm biến. Quay board quanh phương thẳng đứng mà không nghiêng không làm đổi hướng trọng lực tương đối với trục $Z$. Vì thế dữ liệu gia tốc không cung cấp mốc tuyệt đối cho Yaw. Gyroscope trục $Z$ có thể tích phân để theo dõi biến thiên Yaw ngắn hạn, nhưng bias làm kết quả trôi theo thời gian. Phần mở rộng dùng từ kế QMC5883P để tạo mốc phương ngang; đây là một nhánh kỹ thuật riêng vì từ trường chịu ảnh hưởng của kim loại, dòng điện servo và hướng lắp từ kế.

\section{Nguồn sai số trong chuỗi đo}
\subsection{Mô hình từ tín hiệu vật lý đến dữ liệu số}
Một kênh gia tốc có thể được mô tả gần đúng bằng $a^{\rm raw}=s_a a+b_a+n_a$, trong đó $s_a$ là hệ số tỷ lệ, $b_a$ là lệch không và $n_a$ là nhiễu. Kênh gyro tương tự có thể viết $\omega^{\rm raw}=s_g\omega+b_g+n_g$. Sai lệch gá đặt thêm một phép quay giữa hệ trục của board và hệ trục bàn thử. Hiệu chuẩn sáu mặt xử lý thành phần offset và scale của gia tốc; lấy trung bình gyro khi đứng yên xử lý bias tại thời điểm khởi động. Hai bước này chưa xử lý đầy đủ sự thay đổi theo nhiệt độ, sai lệch vuông góc giữa các trục và rung cơ khí.

Nếu sai số bias gyro còn lại là $\delta b$, tích phân trong $T$ giây gây sai số gần $T\delta b$. Ví dụ, chỉ cần bias dư $0{,}02\degree/\mathrm{s}$ thì sau 300~s sai số dự báo đã khoảng $6\degree$. Ví dụ này là hệ quả của phương trình, không phải một số đo mới. Nó giải thích tại sao tích phân gyro cần một mốc dài hạn và tại sao kết quả trôi 330~s ở Chương~3 là phép thử có ý nghĩa.

Sai số gia tốc có dạng khác. Khi board chịu gia tốc tuyến tính, cảm biến đo tổng của trọng lực và gia tốc chuyển động. Thuật toán $\operatorname{atan2}$ không thể tách hai thành phần này từ một mẫu đơn. Rung làm góc gia tốc biến động nhanh, nhất là khi thành phần $a_y$ hoặc $a_z$ thay đổi theo nhịp cơ khí. Bộ lọc bù và Kalman chỉ giảm ảnh hưởng bằng cách đặt trọng số khác nhau cho dự đoán gyro và số đo gia tốc; chúng không biến dữ liệu bị kích động thành góc chuẩn tuyệt đối.

\subsection{Ảnh hưởng của thời gian lấy mẫu}
Phương trình tích phân thường viết với $\Delta t=2\unit{ms}$, tương ứng 500~Hz danh định. Nếu chu kỳ thực tế là $\Delta t_k$, sai số của một bước tích phân xấp xỉ $\omega_k(\Delta t-\Delta t_k)$. Nhiều sai lệch cùng dấu có thể cộng dồn. Vì vậy, cần lưu cả chu kỳ đo thực tế, phân tích trung bình, độ lệch chuẩn, giá trị lớn nhất và số chu kỳ bị trễ. Việc chỉ đọc cấu hình thanh ghi chia tần của MPU6500 không chứng minh toàn bộ hệ thống xử lý liên tục ở 500~Hz.

Với ngắt Data Ready, phần xử lý ngắt ngắn giúp giảm độ trễ phản ứng và bảo vệ các tác vụ khác. Tuy nhiên, một cờ Boolean chỉ thông báo ``có ít nhất một mẫu mới''. Nếu vòng lặp bị giữ quá lâu bởi OLED hoặc I2C, hai ngắt liên tiếp có thể hợp thành một cờ. Lúc ấy tốc độ cập nhật thực của bộ lọc thấp hơn tốc độ sinh mẫu. Đếm số ngắt và số mẫu xử lý trong cùng khoảng thời gian cho phép phát hiện tình huống này. Một phương án cải tiến là lưu dấu thời gian DWT cho từng mẫu và đưa $\Delta t_k$ vào thuật toán sau khi kiểm tra giới hạn hợp lệ.

\subsection{Sai số do chuẩn đối chứng}
Giá trị \,\texttt{ref}\, trong gói telemetry là góc lệnh gửi đến servo, không phải phép đo vị trí trục servo. Servo SG90 có sai khác giữa độ rộng xung và góc thực, có độ rơ khi đổi chiều và có thể không đạt đúng vị trí dưới tải. Sai số đo được theo \,\texttt{ref}\, do đó gộp sai số của thuật toán, cảm biến, gá đặt và servo. Muốn kết luận ``sai số góc Kalman dưới $1\degree$'' cần một chuẩn góc độc lập có bất định đủ nhỏ; có thể dùng thước đo góc đã kiểm tra, encoder hoặc bàn quay được hiệu chuẩn. Nếu chuẩn có bất định gần $1\degree$, kết luận về ngưỡng $1\degree$ sẽ không thuyết phục dù MAE tính toán thấp.

\section{Phân tích chi tiết bốn phương pháp}
\subsection{Góc gia tốc: mốc tuyệt đối có điều kiện}
Công thức ở~\eqref{eq:roll_acc} và~\eqref{eq:pitch_acc} dùng hàm $\operatorname{atan2}$ để xác định góc phần tư và tránh phép chia trực tiếp cho thành phần gần không. Ở tư thế phẳng, khi $a_z$ gần $1g$ và $a_x,a_y$ gần không, cả Roll và Pitch gần $0\degree$. Tại vùng lân cận tư thế đó, một nhiễu nhỏ $\delta a_y$ gây biến thiên Roll xấp xỉ $\delta a_y/a_z$ radian. Nếu board tiến đến tư thế mà vector chiếu xuống mặt phẳng $YZ$ rất nhỏ, góc Roll trở nên nhạy hơn với nhiễu. Đây là một giới hạn hình học, ngoài giới hạn của cảm biến.

Góc gia tốc không tích lũy bias qua thời gian vì mỗi mẫu được tính độc lập. Đổi lại, nó chịu ảnh hưởng trực tiếp của va chạm, rung và gia tốc tuyến tính. Trong cùng bộ dữ liệu, nên luôn giữ đường góc này trên đồ thị để người đọc thấy khi nào phép đo mốc bị nhiễu. Việc bộ lọc cho một đường mượt hơn chỉ chứng minh giảm dao động của tín hiệu đầu ra; độ chính xác thực phải được kiểm tra theo góc chuẩn.

\subsection{Gyro-only: đáp ứng nhanh nhưng thiếu mốc}
Gyro đo vận tốc góc, do đó bắt được biến thiên nhanh ngay cả khi gia tốc kế đang chịu gia tốc tuyến tính. Tích phân rời rạc thuận đơn giản có sai số xấp xỉ phụ thuộc $\Delta t$ và độ biến thiên của vận tốc góc. Một bộ tích phân hình thang có thể giảm sai số rời rạc trong chuyển động nhanh, nhưng không xử lý được bias không đổi. Trong phạm vi phần mềm hiện tại, tích phân thuận được giữ làm đường đối chứng cơ bản để thấy rõ độ trôi và ưu thế của hợp nhất cảm biến.

Giá trị đầu của gyro-only phải được mô tả. Nếu đặt $0\degree$ lúc khởi động trong khi board đang nghiêng, nó thể hiện biến thiên từ thời điểm bật nguồn chứ không phải góc tuyệt đối. Đồ thị động có thể lệch do cả bias và mốc khởi tạo. Khi so sánh bốn phương pháp, cần ghi cùng điều kiện khởi tạo, cùng cửa sổ thời gian và cùng trục, tránh quy lỗi mốc ban đầu thành sai số trôi của thuật toán.

\subsection{Bộ lọc bù dưới góc nhìn đáp ứng tần số}
Phương trình~\eqref{eq:cf} cho thấy gyro đóng vai trò dự đoán ở nhịp nhanh, còn góc gia tốc kéo trạng thái về mốc dài hạn. Xét trường hợp đầu vào gia tốc là hằng số $A$ và gyro bằng không, sai số của bộ lọc sau $n$ bước là $e_n=\alpha^n e_0$. Thời gian cần để sai số còn khoảng $1/e$ giá trị ban đầu là $-\Delta t/\ln\alpha$, gần với $\alpha\Delta t/(1-\alpha)$ khi $\alpha$ gần 1. Với $\alpha=0{,}98$ và $\Delta t=2\unit{ms}$, quy mô thời gian khoảng 0,1~s. Đó là một ước lượng lý thuyết, không phải số đo độ trễ của hệ thống servo.

Nếu $\alpha$ tăng, đường đầu ra ít chịu xung gia tốc hơn nhưng quay về mốc chậm hơn. Nếu $\alpha$ giảm, đáp ứng theo mốc gia tốc nhanh hơn và nhiễu cơ học đi qua nhiều hơn. Tham số thích hợp phụ thuộc độ rung, cấu hình DLPF trong chip và yêu cầu độ trễ. Quét tham số trên một đoạn tĩnh là bước khám phá; cần thêm thử nghiệm STEP có chuẩn góc để chọn tham số theo đồng thời độ ổn định và tốc độ đáp ứng.

\subsection{Kalman hai trạng thái và ý nghĩa của bias}
Mô hình hai trạng thái $[\phi,b]^\mathsf{T}$ cho phép ước lượng cả góc lẫn phần bias gyro còn lại. Dự đoán dùng tốc độ gyro sau khi trừ bias, còn phép đo là góc từ gia tốc. Phần tử $P_{00}$ biểu diễn bất định của góc theo mô hình; $P_{11}$ liên quan bất định của bias. Khi $P_{00}$ lớn hơn tương đối so với $R$, hệ số $K_0=P_{00}^{-}/(P_{00}^{-}+R)$ tăng và bộ lọc tin góc gia tốc nhiều hơn. Khi $R$ lớn, bộ lọc giảm tác động của số đo gia tốc. Hai tham số $Q_{\rm angle}$ và $Q_{\rm bias}$ điều khiển tốc độ tăng bất định khi dự đoán.

Mã \,\texttt{kalman.c}\, hiện dùng $Q_{\rm angle}=0{,}001$, $Q_{\rm bias}=0{,}003$, $R=0{,}03$, với hai đối tượng riêng cho Roll và Pitch. Các giá trị này là tham số kỹ thuật được chọn và thử ngoại tuyến; chúng không phải đặc tính được nhà sản xuất MPU6500 đảm bảo. Vì $P$ khởi tạo bằng không, những mẫu đầu có giai đoạn xác lập. Phép tính sai số tĩnh nên bỏ một khoảng quá độ sau khi bật nguồn và sau mỗi bước chuyển góc, đồng thời nêu rõ độ dài khoảng bỏ.

Việc dùng hai bộ Kalman độc lập giúp chương trình gọn, dễ kiểm tra và đủ cho thử nghiêng từng trục. Nó bỏ qua tương quan giữa Roll và Pitch, sự ghép trục khi quay phức tạp và nhiễu do gia tốc tuyến tính. Trong phép rung mạnh, mô hình đo $z=\phi_{\rm acc}+v$ có thể không còn đúng với giả định nhiễu Gauss gần dừng. Vì thế, mức dao động thấp của Kalman trong một đoạn rung không bảo đảm mọi xung động đều được xử lý tối ưu. Có thể mở rộng bằng kiểm tra độ lớn gia tốc $\lVert\mathbf a\rVert$ trước khi tăng $R$, nhưng phải đo và chứng minh hiệu quả trước khi kết luận.

\section{Biểu diễn số và độ tin cậy của tính toán}
\subsection{Ranh giới của triển khai Q16.16 hiện có}
Ở Q16.16, phép chuyển từ số thực sang trạng thái nguyên dùng hệ số $2^{16}=65\,536$. Với góc trong phạm vi phép thử, độ phân giải danh nghĩa nhỏ hơn nhiều so với nhiễu quan sát. Tuy vậy, độ phân giải không tự động bảo đảm độ chính xác vì còn sai số khi nhân, làm tròn và sai số của cảm biến. Mã \,\texttt{kalman\_fix.c}\, hiện chuyển góc và bias về Q16.16 sau mỗi lần cập nhật; các bước dự đoán hiệp phương sai và tính hệ số Kalman vẫn chạy bằng \,\texttt{float}. Vì vậy, đây là phép thử lượng tử hóa trạng thái, chưa phải cấu trúc Kalman hoàn toàn số nguyên.

Chênh lệch cực đại khoảng $0{,}01\degree$ giữa hai đường float và fixed trên log cũ là một kiểm tra tính nhất quán nội bộ. Nó cho thấy lượng tử hóa trạng thái chưa làm hai đường phân kỳ đáng kể trong bộ dữ liệu đó. Nó không cho biết đường nào gần góc thật hơn. Để so sánh chi phí tính toán, phải đo mỗi hàm riêng, dưới cùng tối ưu compiler, cùng tần số clock và cùng chuỗi dữ liệu vào. Lấy hiệu giữa hai thời gian xử lý tổng ở hai cấu hình khác nhau chỉ cho một ước lượng thô về phần tải bổ sung.

\subsection{Kiểm tra điều kiện biên}
Một phép triển khai đáng tin cậy cần xử lý trường hợp $\Delta t\leq0$, chu kỳ quá lớn do mất mẫu, giá trị cảm biến không hữu hạn và phép tính $S=P_{00}^{-}+R$ không hợp lệ. Trong thiết kế hiện tại, các trạng thái này nên được xem là tình huống cần kiểm thử khi chạy trên mạch cuối. Bảng kiểm thử tối thiểu gồm: giữ yên sau khởi động; quay thuận/ngược cùng biên độ; đổi góc liên tục; mất một mẫu Data Ready; khởi động lại khi bus I2C đang bận; đặt board gần vùng góc lớn. Mỗi trường hợp cần nêu đầu vào, phản ứng mong đợi và log xác nhận.

Hiệu chuẩn gyro cũng phải được lặp lại khi thay nhiệt độ hoặc sau khi gá mạch khác. Nếu hiệu chuẩn diễn ra trong lúc servo chưa đứng yên, trung bình 1\,000 mẫu sẽ tính chuyển động thành bias và gây sai số ở mọi phép thử sau đó. Quy trình khởi động nên giữ servo ở vị trí cố định, chờ dao động cơ học tắt rồi mới lấy bias. Thứ tự khởi tạo này là điều kiện thực nghiệm, không chỉ là chi tiết giao diện phần mềm.
